\section{Declaración de uso de inteligencia artificial}\label{sec:ia}

\textit{Herramienta y versión.} Claude, de Anthropic, modelo Claude Opus 5.5, usado desde Claude Code en la aplicación de escritorio para Windows, entre el 20 y el 22 de noviembre de 2026.

\textit{Para qué se usó.} El equipo usó Claude Code en todas las secciones de este documento, en el esqueleto, en los ADR, en los diagramas y en la búsqueda de referencias:
\begin{itemize}
  \item Borradores de todas las secciones del documento y de los ADR, y la conversión del material a \LaTeX. El resumen, el contexto, los objetivos y los requisitos funcionales se redactaron a partir del documento de requisitos de Aplicaciones Web del equipo y del Reglamento de la UIC; las historias HU-01 a HU-08 se copiaron de ese documento. Los riesgos R-08 a R-10 se redactaron a partir de las lecciones que el equipo ya había escrito en FORM-AC-01.
  \item Código del esqueleto (ms-usuarios y ms-propuestas), sus pruebas automáticas y el script que crea los issues.
  \item Generación de los diagramas C4 en formato draw.io y su exportación a las figuras del documento.
  \item Búsqueda de referencias y consulta de sus metadatos en Crossref.
\end{itemize}

\textit{Para qué no se usó.} La entrevista con la coordinación de carrera, la consulta a docentes, las fotografías y la grabación son trabajo del equipo. Claude Code no aportó datos de la entrevista: el documento no atribuye frases a la entrevistada ni a los docentes.

\textit{Cómo lo verificó el equipo.}
\begin{itemize}
  \item El código se compiló y pasó sus 9 pruebas automáticas; luego se ejecutaron los dos servicios y se comprobaron la respuesta combinada, el 404 y el 503 con ms-usuarios apagado. Las capturas de la sección~\ref{sec:esqueleto} son de esa ejecución.
  \item Cada referencia con DOI se comparó con lo que registra Crossref; se descartaron cinco que no coincidían. El informe de Meta se comprobó en su página (apéndice~\ref{sec:verificacion}).
  \item Los diagramas se abrieron con el visor de draw.io para revisar que cada flecha coincida con la tabla~\ref{tab:interacciones}.
  \item En el contexto, los objetivos y los requisitos, cada dato indica la sección del documento de Aplicaciones Web o el artículo del reglamento de donde sale. Lo que no tenía fuente se redactó sin cifras.
  \item El documento se compiló con la secuencia del README (pdflatex, biber, pdflatex, pdflatex) sin errores.
  \item Cada responsable revisó su parte antes de subirla al repositorio y debe poder explicarla sin ayuda.
\end{itemize}
